% GENERATED FILE. Edit canonical lessons, then run npm run generate:books.
% canonical-source-hash: fnv1a64:6e26b00936c3db8f
% canonical-lessons: AR-C433-itar, AR-C433-ajala, AR-C433-thuqb, AR-C433-midakhkha, AR-C433-jarra

\chapter{A flat tyre}
\label{ch:ar-a-flat-tyre}
\hlchaptermodality{433}

\begin{chapteropening}
\textbf{By the end of this chapter:} \emph{I can say a tyre, a wheel, a hole, a pump and to tow.}

Complete the last lesson of chapter 433: I can say a tyre, a wheel, a hole, a pump and to tow.
\end{chapteropening}

\section[\texorpdfstring{\ar{إطار} (iṭār)}{إطار (iṭār)}]{\ar{إطار} (iṭār) — a tyre; a frame}

\label{lesson:AR-C433-itar}



\begin{quote}
Before the new one: say the Arabic for brakes, then the Arabic for to supply.
\end{quote}

\subsection*{You'll want to know: \ar{إطار}}
\textbf{\ar{إطار}} — \emph{iṭār} — \textquotedblleft{}a tyre; a frame\textquotedblright{}.

\begin{quote}
\ar{كان} \ar{إطار} \ar{الدراجة} \ar{قديما}.
\end{quote}

\begin{quote}
\emph{kāna iṭār ad-darrāja qadīman.}
\end{quote}

\begin{quote}
The bicycle's tyre was old.
\end{quote}

\begin{cousinweb}
\textbf{\ar{إطار}} (\emph{iṭār}) is usually traced to the native Arabic root \emph{ʾ-ṭ-r}, \textquotedblleft{}to bend round\textquotedblright{}. An \emph{iṭār} is a ring or rim that goes round a thing.
\end{cousinweb}

\begin{grammarlens}[title={What you've built}]
One more word for getting around.
\end{grammarlens}

\subsection*{Your turn}
\begin{itemize}
\raggedright
  \item \emph{Say it:} \emph{iṭār}
  \item \emph{Say it:} \emph{iṭār}, once more
  \item \emph{Say it:} \emph{kāna iṭār ad-darrāja qadīman.} — The bicycle's tyre was old.
  \item \emph{Recall:} say the Arabic for brakes, then the Arabic for to supply, then say \emph{iṭār} again
\end{itemize}

\begin{tcolorbox}[breakable,colback=teal!4,colframe=teal!35!black,title={Before you move on}]
What does \ar{إطار} mean? (\textquotedblleft{}A tyre; a frame\textquotedblright{}.) Say it once more.
\end{tcolorbox}
\section[\texorpdfstring{\ar{عجلة} (ʿajala)}{عجلة (ʿajala)}]{\ar{عجلة} (ʿajala) — a wheel}

\label{lesson:AR-C433-ajala}



\begin{quote}
Before the new one: say the Arabic for to supply, then the Arabic for a tyre.
\end{quote}

\subsection*{You'll want to know: \ar{عجلة}}
\textbf{\ar{عجلة}} — \emph{ʿajala} — \textquotedblleft{}a wheel\textquotedblright{}.

\begin{quote}
\ar{غير} \ar{أبي} \ar{عجلة} \ar{السيارة} \ar{على} \ar{الطريق}.
\end{quote}

\begin{quote}
\emph{ghayyara abī ʿajalat as-sayyāra ʿalā aṭ-ṭarīq.}
\end{quote}

\begin{quote}
My father changed the car's wheel on the road.
\end{quote}

The same word means haste: \emph{fī ʿajala}, \textquotedblleft{}in a hurry\textquotedblright{}.

\begin{cousinweb}
\textbf{\ar{عجلة}} (\emph{ʿajala}) is built on the native Arabic root \emph{ʿ-j-l}, \textquotedblleft{}to hasten\textquotedblright{} (Lane). It is the root of \emph{mustaʿjil}, \textquotedblleft{}in a hurry\textquotedblright{}; Lane lists \emph{ʿajala} as a cart drawn by an ox, said to be named for its quick going, and modern Arabic uses it for a wheel.
\end{cousinweb}

\begin{grammarlens}[title={What you've built}]
One more word for getting around.
\end{grammarlens}

\subsection*{Your turn}
\begin{itemize}
\raggedright
  \item \emph{Say it:} \emph{ʿajala}
  \item \emph{Say it:} \emph{ʿajala}, once more
  \item \emph{Say it:} \emph{ghayyara abī ʿajalat as-sayyāra ʿalā aṭ-ṭarīq.} — My father changed the car's wheel on the road.
  \item \emph{Recall:} say the Arabic for to supply, then the Arabic for a tyre, then say \emph{ʿajala} again
\end{itemize}

\begin{tcolorbox}[breakable,colback=teal!4,colframe=teal!35!black,title={Before you move on}]
What does \ar{عجلة} mean? (\textquotedblleft{}A wheel\textquotedblright{}.) Say it once more.
\end{tcolorbox}
\section[\texorpdfstring{\ar{ثقب} (thuqb)}{ثقب (thuqb)}]{\ar{ثقب} (thuqb) — a hole; a puncture}

\label{lesson:AR-C433-thuqb}



\begin{quote}
Before the new one: say the Arabic for a tyre, then the Arabic for a wheel.
\end{quote}

\subsection*{You'll want to know: \ar{ثقب}}
\textbf{\ar{ثقب}} — \emph{thuqb} — \textquotedblleft{}a hole; a puncture\textquotedblright{}.

\begin{quote}
\ar{وجد} \ar{العامل} \ar{ثقبا} \ar{صغيرا} \ar{في} \ar{الإطار}.
\end{quote}

\begin{quote}
\emph{wajada al-ʿāmil thuqban ṣaghīran fī al-iṭār.}
\end{quote}

\begin{quote}
The worker found a small hole in the tyre.
\end{quote}

\begin{cousinweb}
\textbf{\ar{ثقب}} (\emph{thuqb}) is built on the native Arabic root \emph{th-q-b}, \textquotedblleft{}to pierce\textquotedblright{} (Lane).
\end{cousinweb}

\begin{grammarlens}[title={What you've built}]
One more word for getting around.
\end{grammarlens}

\subsection*{Your turn}
\begin{itemize}
\raggedright
  \item \emph{Say it:} \emph{thuqb}
  \item \emph{Say it:} \emph{thuqb}, once more
  \item \emph{Say it:} \emph{wajada al-ʿāmil thuqban ṣaghīran fī al-iṭār.} — The worker found a small hole in the tyre.
  \item \emph{Recall:} say the Arabic for a tyre, then the Arabic for a wheel, then say \emph{thuqb} again
\end{itemize}

\begin{tcolorbox}[breakable,colback=teal!4,colframe=teal!35!black,title={Before you move on}]
What does \ar{ثقب} mean? (\textquotedblleft{}A hole; a puncture\textquotedblright{}.) Say it once more.
\end{tcolorbox}
\section[\texorpdfstring{\ar{مضخة} (miḍakhkha)}{مضخة (miḍakhkha)}]{\ar{مضخة} (miḍakhkha) — a pump}

\label{lesson:AR-C433-midakhkha}



\begin{quote}
Before the new one: say the Arabic for a wheel, then the Arabic for a hole.
\end{quote}

\subsection*{You'll want to know: \ar{مضخة}}
\textbf{\ar{مضخة}} — \emph{miḍakhkha} — \textquotedblleft{}a pump\textquotedblright{}.

\begin{quote}
\ar{نفخ} \ar{أخي} \ar{الإطار} \ar{بمضخة} \ar{صغيرة}.
\end{quote}

\begin{quote}
\emph{nafakha akhī al-iṭār bi-miḍakhkha ṣaghīra.}
\end{quote}

\begin{quote}
My brother pumped up the tyre with a small pump.
\end{quote}

The pattern \emph{mifʿala} names a tool.

\begin{cousinweb}
\textbf{\ar{مضخة}} (\emph{miḍakhkha}) is usually traced to the native Arabic root \emph{ḍ-kh-kh}, \textquotedblleft{}to spurt (of water)\textquotedblright{}.
\end{cousinweb}

\begin{grammarlens}[title={What you've built}]
One more word for getting around.
\end{grammarlens}

\subsection*{Your turn}
\begin{itemize}
\raggedright
  \item \emph{Say it:} \emph{miḍakhkha}
  \item \emph{Say it:} \emph{miḍakhkha}, once more
  \item \emph{Say it:} \emph{nafakha akhī al-iṭār bi-miḍakhkha ṣaghīra.} — My brother pumped up the tyre with a small pump.
  \item \emph{Recall:} say the Arabic for a wheel, then the Arabic for a hole, then say \emph{miḍakhkha} again
\end{itemize}

\begin{tcolorbox}[breakable,colback=teal!4,colframe=teal!35!black,title={Before you move on}]
What does \ar{مضخة} mean? (\textquotedblleft{}A pump\textquotedblright{}.) Say it once more.
\end{tcolorbox}
\section[\texorpdfstring{\ar{جر} (jarra)}{جر (jarra)}]{\ar{جر} (jarra) — to tow; to drag}

\label{lesson:AR-C433-jarra}



\begin{quote}
Before the new one: say the Arabic for a hole, then the Arabic for a pump.
\end{quote}

\subsection*{You'll want to know: \ar{جر}}
\textbf{\ar{جر}} — \emph{jarra} — \textquotedblleft{}to tow; to drag\textquotedblright{}.

\begin{quote}
\ar{جر} \ar{رجل} \ar{سيارتنا} \ar{إلى} \ar{القرية}.
\end{quote}

\begin{quote}
\emph{jarra rajul sayyāratanā ilā al-qarya.}
\end{quote}

\begin{quote}
A man towed our car to the village.
\end{quote}

Present: \ar{يجر}, \emph{yajurr}.

Grammarians call the genitive case \emph{al-jarr}, \textquotedblleft{}the dragging\textquotedblright{}.

\begin{cousinweb}
\textbf{\ar{جر}} (\emph{jarra}) is built on the native Arabic root \emph{j-r-r}, \textquotedblleft{}to drag, to pull\textquotedblright{} (Lane). It is the plain verb of a doubled root.
\end{cousinweb}

\begin{grammarlens}[title={What you've built}]
One more word for getting around.
\end{grammarlens}

\subsection*{Your turn}
\begin{itemize}
\raggedright
  \item \emph{Say it:} \emph{jarra}
  \item \emph{Say it:} \emph{jarra}, once more
  \item \emph{Say it:} \emph{jarra rajul sayyāratanā ilā al-qarya.} — A man towed our car to the village.
  \item \emph{Recall:} say the Arabic for a hole, then the Arabic for a pump, then say \emph{jarra} again
\end{itemize}

\begin{tcolorbox}[breakable,colback=teal!4,colframe=teal!35!black,title={Before you move on}]
What does \ar{جر} mean? (\textquotedblleft{}To tow; to drag\textquotedblright{}.) Say it once more.
\end{tcolorbox}
